\scanpage{185}\mspage{41}
% archive sheet number: 185

La notion de $W_0$-fibration apparaît comme extrêmement forte -- beaucoup plus forte à certains égards que celle (\unsure{disons|absolue}) de $W_\infty$-fibration \inserted{(\unsure{du moins|au moins} \unsure{ostensiblement})} \deleted{$W_0$-}fibration. Elle \unsure{équivaut|revient} à ceci : pour tte composante connexe $I'_\alpha$ de $I'_\alpha$, désignant par $I_\alpha$ la composante connexe correspondante de $I$, et pour tout $J_\alpha$ $0$-connexe \inserted{et} au dessus de $I_\alpha$, le produit fibré $I'_\alpha\times_{I_\alpha}J_\alpha$ est encore $0$-connexe. (\unsure{Ça|Sa} ressemble aux histoires de totale $0$-connexité d'un topos !) \deleted{Cel\unsure{a} implique que les fibres de $I'_\alpha\longrightarrow I_\alpha$ sont $0$-connexes (prendre $J_\alpha=e_{\unsure{i}}$), et \unsure{a}} Mais cela signifie \deleted{\unsure{que|pour}} que $I'_\alpha\longrightarrow I_\alpha$ est \inserted{une $W_0$-équivalence} universelle\deleted{\unsure{ment}}. Cela implique que ses fibres sont $0$-connexes, et la réciproque est vraie si $u:I'\longrightarrow I$ est $W_0$-propre ou $W_0$-lisse. Ainsi, la condition \unsure{essentielle|en italiques}$\unsure{\Rightarrow|,}$ pour une \inserted{$W_0$-fibration} est celle que les fibres des flèches induites $I'_\alpha\longrightarrow I_{\unsure{\beta(\alpha)|\mu(\alpha)}}$ soient $0$-connexes ; \deleted{\unsure{si on suppose} \unsure{2\textsuperscript{o}} que pour une flèche de $I$ i.e.} \inserted{En fait}

\uuline{\deleted{\unsure{proposition}} Proposition 3.1} Soit $u:I'\longrightarrow I$ une flèche dans $\mathrm{Cat}$.

a) Pour que $u$ soit une $W_0$-équivalence universelle, il f et s. que \deleted{ses fibres soient $W_0$-connexes,} pour toute \deleted{flèche $\alpha$ de $I$} $\Delta^1\longrightarrow I$ i.e. \inserted{toute flèche $\alpha$ de $I$}, \deleted{et il suffit} la catégorie $I'\times_I\Delta^1\overset{\mathrm{déf}}{=}\deleted{\unsure{I}}I'|\alpha$ soit \deleted{\unsure{les points} $\alpha$ \unsure{mis} \unsure{bout}} $0$-connexe, ou encore que les fibres de $I$ soient $0$-connexes, et que pour toute flèche \uline{non identique} $\alpha$ de $I$, $I'|\alpha$ soit $0$-connexe.

b) Pour que $u:\deleted{\unsure{V}}I'\longrightarrow I$ soit une $W_0$-fibration, il f et s. que pour tte composante connexe $I'_\alpha$, de
